\begin{enumerate}[label=(\alph*)]
\item Do năng lượng bảo toàn nên:
\begin{equation}
    qA\Delta t = L_i\rho_idA.
\end{equation}
\begin{equation}
   \Longrightarrow d=\frac{q\Delta t}{L_i\rho_i}=6.1 \si{mm}.
\end{equation}
\item 
    \begin{enumerate}[label=(\roman*)]
    \item Ta đặt hệ trục tọa độ như hình vẽ bên dưới với $y_1$, $y_2$ lần lượt biểu diễn hình dạng của mặt dưới và mặt trên của lớp băng, ta tính áp suất tại một điểm nằm ngay dưới lớp băng:

    \begin{figure}[H]
    \centering



% Pattern Info
 
\tikzset{
pattern size/.store in=\mcSize, 
pattern size = 5pt,
pattern thickness/.store in=\mcThickness, 
pattern thickness = 0.3pt,
pattern radius/.store in=\mcRadius, 
pattern radius = 1pt}
\makeatletter
\pgfutil@ifundefined{pgf@pattern@name@_bbnil2yq2}{
\pgfdeclarepatternformonly[\mcThickness,\mcSize]{_bbnil2yq2}
{\pgfqpoint{0pt}{0pt}}
{\pgfpoint{\mcSize+\mcThickness}{\mcSize+\mcThickness}}
{\pgfpoint{\mcSize}{\mcSize}}
{
\pgfsetcolor{\tikz@pattern@color}
\pgfsetlinewidth{\mcThickness}
\pgfpathmoveto{\pgfqpoint{0pt}{0pt}}
\pgfpathlineto{\pgfpoint{\mcSize+\mcThickness}{\mcSize+\mcThickness}}
\pgfusepath{stroke}
}}
\makeatother

% Pattern Info
 
\tikzset{
pattern size/.store in=\mcSize, 
pattern size = 5pt,
pattern thickness/.store in=\mcThickness, 
pattern thickness = 0.3pt,
pattern radius/.store in=\mcRadius, 
pattern radius = 1pt}
\makeatletter
\pgfutil@ifundefined{pgf@pattern@name@_k8wwxhed7}{
\pgfdeclarepatternformonly[\mcThickness,\mcSize]{_k8wwxhed7}
{\pgfqpoint{0pt}{0pt}}
{\pgfpoint{\mcSize+\mcThickness}{\mcSize+\mcThickness}}
{\pgfpoint{\mcSize}{\mcSize}}
{
\pgfsetcolor{\tikz@pattern@color}
\pgfsetlinewidth{\mcThickness}
\pgfpathmoveto{\pgfqpoint{0pt}{0pt}}
\pgfpathlineto{\pgfpoint{\mcSize+\mcThickness}{\mcSize+\mcThickness}}
\pgfusepath{stroke}
}}
\makeatother

% Pattern Info
 
\tikzset{
pattern size/.store in=\mcSize, 
pattern size = 5pt,
pattern thickness/.store in=\mcThickness, 
pattern thickness = 0.3pt,
pattern radius/.store in=\mcRadius, 
pattern radius = 1pt}
\makeatletter
\pgfutil@ifundefined{pgf@pattern@name@_tidfkeuug}{
\pgfdeclarepatternformonly[\mcThickness,\mcSize]{_tidfkeuug}
{\pgfqpoint{0pt}{0pt}}
{\pgfpoint{\mcSize}{\mcSize}}
{\pgfpoint{\mcSize}{\mcSize}}
{
\pgfsetcolor{\tikz@pattern@color}
\pgfsetlinewidth{\mcThickness}
\pgfpathmoveto{\pgfqpoint{0pt}{\mcSize}}
\pgfpathlineto{\pgfpoint{\mcSize+\mcThickness}{-\mcThickness}}
\pgfpathmoveto{\pgfqpoint{0pt}{0pt}}
\pgfpathlineto{\pgfpoint{\mcSize+\mcThickness}{\mcSize+\mcThickness}}
\pgfusepath{stroke}
}}
\makeatother

% Pattern Info
 
\tikzset{
pattern size/.store in=\mcSize, 
pattern size = 5pt,
pattern thickness/.store in=\mcThickness, 
pattern thickness = 0.3pt,
pattern radius/.store in=\mcRadius, 
pattern radius = 1pt}
\makeatletter
\pgfutil@ifundefined{pgf@pattern@name@_lv9hs0y0o}{
\pgfdeclarepatternformonly[\mcThickness,\mcSize]{_lv9hs0y0o}
{\pgfqpoint{0pt}{0pt}}
{\pgfpoint{\mcSize}{\mcSize}}
{\pgfpoint{\mcSize}{\mcSize}}
{
\pgfsetcolor{\tikz@pattern@color}
\pgfsetlinewidth{\mcThickness}
\pgfpathmoveto{\pgfqpoint{0pt}{\mcSize}}
\pgfpathlineto{\pgfpoint{\mcSize+\mcThickness}{-\mcThickness}}
\pgfpathmoveto{\pgfqpoint{0pt}{0pt}}
\pgfpathlineto{\pgfpoint{\mcSize+\mcThickness}{\mcSize+\mcThickness}}
\pgfusepath{stroke}
}}
\makeatother
\tikzset{every picture/.style={line width=0.75pt}} %set default line width to 0.75pt        

\begin{tikzpicture}[x=0.75pt,y=0.75pt,yscale=-1,xscale=1]
%uncomment if require: \path (0,413); %set diagram left start at 0, and has height of 413

%Straight Lines [id:da5230596040478203] 
\draw    (37.47,366.5) -- (462.92,366.5) ;
\draw [shift={(464.92,366.5)}, rotate = 180] [color={rgb, 255:red, 0; green, 0; blue, 0 }  ][line width=0.75]    (10.93,-3.29) .. controls (6.95,-1.4) and (3.31,-0.3) .. (0,0) .. controls (3.31,0.3) and (6.95,1.4) .. (10.93,3.29)   ;
%Straight Lines [id:da7662024978089438] 
\draw    (37.47,366.5) -- (37.47,37.5) ;
\draw [shift={(37.47,35.5)}, rotate = 90] [color={rgb, 255:red, 0; green, 0; blue, 0 }  ][line width=0.75]    (10.93,-3.29) .. controls (6.95,-1.4) and (3.31,-0.3) .. (0,0) .. controls (3.31,0.3) and (6.95,1.4) .. (10.93,3.29)   ;
%Straight Lines [id:da8418157269690355] 
\draw    (38.71,184.83) -- (310.05,52.88) ;
%Curve Lines [id:da17007007313612144] 
\draw [color={rgb, 255:red, 0; green, 0; blue, 0 }  ,draw opacity=1 ][pattern=_bbnil2yq2,pattern size=6pt,pattern thickness=0.75pt,pattern radius=0pt, pattern color={rgb, 255:red, 0; green, 0; blue, 0}]   (37.47,366.5) .. controls (186.15,211.63) and (367.04,233.93) .. (401.73,247.56) ;
%Straight Lines [id:da04869161967510438] 
\draw  [dash pattern={on 3.75pt off 3pt on 7.5pt off 1.5pt}]  (38.71,184.83) -- (329,43.5) ;
%Shape: Right Triangle [id:dp9262687264046802] 
\draw  [draw opacity=0][pattern=_k8wwxhed7,pattern size=6pt,pattern thickness=0.75pt,pattern radius=0pt, pattern color={rgb, 255:red, 0; green, 0; blue, 0}] (401.73,246.32) -- (37.47,366.5) -- (401.73,366.5) -- cycle ;
%Curve Lines [id:da2863251511770848] 
\draw    (38.71,349.15) .. controls (187.39,194.28) and (368.28,216.58) .. (402.97,230.21) ;
%Shape: Parallelogram [id:dp8689460533005648] 
\draw  [pattern=_tidfkeuug,pattern size=6pt,pattern thickness=0.75pt,pattern radius=0pt, pattern color={rgb, 255:red, 0; green, 0; blue, 0}] (229.57,233.88) -- (248.84,229.19) -- (248.27,247.8) -- (229,252.5) -- cycle ;
%Straight Lines [id:da028476772170188402] 
\draw  [dash pattern={on 3.75pt off 3pt on 7.5pt off 1.5pt}]  (229.11,246.56) -- (225,188.5) ;
%Straight Lines [id:da5851897809830424] 
\draw  [dash pattern={on 3.75pt off 3pt on 7.5pt off 1.5pt}]  (248.84,229.19) -- (248.84,186.06) ;
%Straight Lines [id:da03823942730739993] 
\draw    (225,182.38) -- (248.2,180.96) ;
\draw [shift={(250.19,180.84)}, rotate = 176.5] [color={rgb, 255:red, 0; green, 0; blue, 0 }  ][line width=0.75]    (6.56,-1.97) .. controls (4.17,-0.84) and (1.99,-0.18) .. (0,0) .. controls (1.99,0.18) and (4.17,0.84) .. (6.56,1.97)   ;
\draw [shift={(223,182.5)}, rotate = 356.5] [color={rgb, 255:red, 0; green, 0; blue, 0 }  ][line width=0.75]    (6.56,-1.97) .. controls (4.17,-0.84) and (1.99,-0.18) .. (0,0) .. controls (1.99,0.18) and (4.17,0.84) .. (6.56,1.97)   ;
%Straight Lines [id:da4179548287061784] 
\draw  [dash pattern={on 3.75pt off 3pt on 7.5pt off 1.5pt}]  (38,240.86) -- (237.3,240.86) ;
%Shape: Parallelogram [id:dp557566449587727] 
\draw  [pattern=_lv9hs0y0o,pattern size=6pt,pattern thickness=0.75pt,pattern radius=0pt, pattern color={rgb, 255:red, 0; green, 0; blue, 0}] (128.8,277.43) -- (149.66,266.63) -- (153.85,279.7) -- (133,290.5) -- cycle ;
%Straight Lines [id:da7633202047345914] 
\draw  [dash pattern={on 3.75pt off 3pt on 7.5pt off 1.5pt}]  (37,278.56) -- (141.33,278.56) ;
%Straight Lines [id:da3552047096139397] 
\draw  [dash pattern={on 3.75pt off 3pt on 7.5pt off 1.5pt}]  (133,290.5) -- (114,221.5) ;
%Straight Lines [id:da13298156855753085] 
\draw  [dash pattern={on 3.75pt off 3pt on 7.5pt off 1.5pt}]  (153.85,279.7) -- (138,217.5) ;
%Straight Lines [id:da4898700068171965] 
\draw    (115.86,220.76) -- (132.14,214.24) ;
\draw [shift={(134,213.5)}, rotate = 158.2] [color={rgb, 255:red, 0; green, 0; blue, 0 }  ][line width=0.75]    (6.56,-1.97) .. controls (4.17,-0.84) and (1.99,-0.18) .. (0,0) .. controls (1.99,0.18) and (4.17,0.84) .. (6.56,1.97)   ;
\draw [shift={(114,221.5)}, rotate = 338.2] [color={rgb, 255:red, 0; green, 0; blue, 0 }  ][line width=0.75]    (6.56,-1.97) .. controls (4.17,-0.84) and (1.99,-0.18) .. (0,0) .. controls (1.99,0.18) and (4.17,0.84) .. (6.56,1.97)   ;
%Straight Lines [id:da9777957517018839] 
\draw    (95,306.5) -- (127.3,286.55) ;
\draw [shift={(129,285.5)}, rotate = 148.3] [color={rgb, 255:red, 0; green, 0; blue, 0 }  ][line width=0.75]    (10.93,-3.29) .. controls (6.95,-1.4) and (3.31,-0.3) .. (0,0) .. controls (3.31,0.3) and (6.95,1.4) .. (10.93,3.29)   ;
%Straight Lines [id:da16384066392461127] 
\draw    (288,232.5) -- (251.98,237.24) ;
\draw [shift={(250,237.5)}, rotate = 352.5] [color={rgb, 255:red, 0; green, 0; blue, 0 }  ][line width=0.75]    (10.93,-3.29) .. controls (6.95,-1.4) and (3.31,-0.3) .. (0,0) .. controls (3.31,0.3) and (6.95,1.4) .. (10.93,3.29)   ;
%Shape: Rectangle [id:dp40868365246318505] 
\draw  [fill={rgb, 255:red, 255; green, 255; blue, 255 }  ,fill opacity=1 ] (111,304) -- (149,304) -- (149,326.5) -- (111,326.5) -- cycle ;
%Shape: Rectangle [id:dp6853395673450766] 
\draw  [fill={rgb, 255:red, 255; green, 255; blue, 255 }  ,fill opacity=1 ] (255,248) -- (298,248) -- (298,270.5) -- (255,270.5) -- cycle ;

% Text Node
\draw (35.47,32.1) node [anchor=south east] [inner sep=0.75pt]    {$y$};
% Text Node
\draw (468.83,360.47) node [anchor=south west] [inner sep=0.75pt]    {$x$};
% Text Node
\draw (291.22,127.63) node [anchor=north west][inner sep=0.75pt]   [align=left] {{\fontfamily{lmss}\selectfont Băng}};
% Text Node
\draw (325.56,47.18) node [anchor=north west][inner sep=0.75pt]    {$y_{2}$};
% Text Node
\draw (404.97,233.61) node [anchor=north west][inner sep=0.75pt]    {$y_{1}$};
% Text Node
\draw (225.31,157.88) node [anchor=north west][inner sep=0.75pt]    {$\Delta l'$};
% Text Node
\draw (36,240.86) node [anchor=east] [inner sep=0.75pt]    {$y_{1} '$};
% Text Node
\draw (293.03,225) node [anchor=north west][inner sep=0.75pt]   [align=left] {{\fontfamily{lmss}\selectfont Nước}};
% Text Node
\draw (35,278.56) node [anchor=east] [inner sep=0.75pt]    {$y_{1}$};
% Text Node
\draw (110.31,190.88) node [anchor=north west][inner sep=0.75pt]    {$\Delta l$};
% Text Node
\draw (113,307.4) node [anchor=north west][inner sep=0.75pt]    {$p.A$};
% Text Node
\draw (257,251.4) node [anchor=north west][inner sep=0.75pt]    {$p'.A'$};


\end{tikzpicture}
\caption{Mặt cắt của chỏm băng với nước ở giữa chỏm băng với đất}
    \label{fig:S5_01}
    \end{figure}

    Với $p_a$ là áp suất khí quyển, ta dễ dàng thấy được áp suất tại một điểm ngay dưới lớp băng là: 
    \begin{equation}
        p=\rho_ig(y_2-y_1) + p_a.
    \end{equation}
    Bây giờ ta xét một ống nước từ $y_1$ đến $y_1'$, có diện tích $A$, $A'$, áp suất tại đó là $p$, $p'$. Xét một sự dịch chuyển ảo nhỏ sao cho đầu $y_1$ và $y_1'$ dịch đi một lượng nhỏ là $\Delta l$, $\Delta l'$. Khi đó công do áp lực tác dụng lên dòng chất lỏng:
    \begin{equation}
        \Delta W_1 = - p'A'\Delta l' - (-p A\Delta l ) = p A\Delta l - p'A'\Delta l'.
    \end{equation}
    Công do trọng lực tác dụng:
    \begin{equation}
        \Delta W_2= -\rho_w\Delta l'A'gy_1'-(-\rho_w\Delta lAgy_1)=\rho_w\Delta lAgy_1-\rho_w\Delta l'A'gy_1'.
    \end{equation}

    Theo nguyên lí công ảo: $\Delta W_1 + \Delta W_2=0$:
    \begin{equation}
    \label{eq:S5_06}
        \Longrightarrow p A\Delta l - p'A'\Delta l' + \rho_w\Delta lAgy_1-\rho_w\Delta l'A'gy_1'=0.
    \end{equation}
    Do khối lượng của dòng nước bảo toàn, tức là toàn bộ lượng khối lượng đi vào ($\rho_wA\Delta l$) phải bằng khối lượng đi ra ($\rho_wA'\Delta l'$), hay:
    \begin{equation}
        A\Delta l=A'\Delta l'.
    \end{equation}
    Thay vào phương trình \ref{eq:S5_06} ta thu được:
    \begin{equation}
        p + \rho_wgy_1=p'+\rho_wgy_1'.
    \end{equation}
    \begin{equation}
    \label{eq:S5_10}
        \Longrightarrow p + \rho_wgy_1=const.
    \end{equation}
    \begin{equation}
        \Longrightarrow \rho_ig( y_2- y_1) + \rho_w y_1 g= const.
    \end{equation}

    Ta có thể chỉ ra điều này dễ dàng hơn thông qua phương trình Bernoulli.
    
    Do băng chỉ chuyển động thẳng đứng, ta có thể chỉ ra rằng bán kính hình nón này cũng bằng bán kính $r$ của khu vực bị trũng xuống ở bề mặt chỏm băng.
    
    Với $y_2= \frac{h}{r}x + D$:
    \begin{equation}
        \Longrightarrow y_1= -\frac{\rho_i}{\rho_w-\rho_i}\frac{h}{r}x +\frac{\rho_i}{\rho_w-\rho_i}h.
    \end{equation}
    Từ đây ta có thể thấy $y_1$ có dạng đường thẳng. Ngoài ra, do hệ số góc của $y_1$ âm và do đối xứng quay, phần nước ở dưới chỏm băng có hình dạng là một hình nón úp. Ta dễ thấy liên hệ độ sâu khu vực bị trũng xuống với chiều cao hình nón được tạo thành:
    \begin{equation}
    \label{eq:S5_12}
        H=\frac{\rho_i}{\rho_w-\rho_i}h.
    \end{equation}
    \item Đầu tiên, khi magma mới thâm nhập vào trong băng, trạng thái cân bằng thủy tĩnh chưa được thiết lập nên toàn bộ khối lượng băng có hình nón bán kính $r$, chiều cao $h_1$ tan hết thành nước và đều bị chảy đi hết, bề mặt chỏm băng vẫn nằm ngang. Sau đó trạng thái cân bằng thủy tĩnh được thiết lập nhanh chóng, nước chảy ra ngoài thêm khiến cho bề mặt băng bị trũng xuống một lượng $h_2=\frac{\rho_w-\rho_i}{\rho_i}h_1$ (Sử dụng biểu thức \ref{eq:S5_12}). Sau khi trạng thái cân bằng thủy tĩnh được thiết lập, khối magma nguội đi khiến cho băng chảy thêm nhưng mà không bị chảy nước ra ngoài nữa mà thay vào đó chiếm chỗ ở phần bên trên magma và làm cho băng bị trũng xuống thêm. Gọi chiều cao của hình nón băng tan chảy đó là $h_3$, khi đó phần chiếm chỗ có thể tích tương đương hình nón có chiều cao $h_3'=\frac{\rho_i}{\rho_w}h_3$ do không có lượng nước nào chảy ra ngoài nữa, bề mặt băng khi đó bị trũng xuống thêm đoạn $h_4=\frac{\rho_w-\rho_i}{\rho_i}h_3'$.

    Độ sâu của khu vực bị trũng là: 
    \begin{equation}
        h=\frac{\rho_w-\rho_i}{\rho_i}(h_1+h_3') = \frac{\rho_w-\rho_i}{\rho_i}H.
    \end{equation}
    \begin{equation}
        \Longrightarrow H=\frac{\rho_i}{\rho_w-\rho_i}h = 1105\si{m}.
    \end{equation}

    Ta có tổng khối lượng băng bị tan:
    \begin{equation}
        m_{tot}= \rho_i\frac{1}{3}\pi r^2(h_1+h_2+h_3).
    \end{equation}   

    Phương trình cân bằng nhiệt:
    \begin{equation}
    \label{eq:S5_16}
        L_r\rho_r\frac{1}{3}\pi r^2h_1 + \rho_r\frac{1}{3}\pi r^2h_1c_r\Delta T= L_im_{tot} = L_i\rho_i\frac{1}{3}\pi r^2(h_1+h_2+h_3).
    \end{equation}

    Thay $h_2=\frac{\rho_w-\rho_i}{\rho_i}h_1$, $H-h_1=h_3'=\frac{\rho_i}{\rho_w}h_3$ vào phương trình \ref{eq:S5_16}, ta thu được:
    \begin{equation}
        h_1=\frac{\rho_i\rho_wL_ih}{(\rho_w-\rho_i)(L_r+c_r\Delta T)}= 103\si{m}.
    \end{equation}
    Từ đó ta dễ tính được khối lượng băng đã chảy:
    \begin{equation}
        m_{tot}= \rho_i\frac{1}{3}\pi r^2(h_1+h_2+h_3) = 2.9\times10^{11}\si{kg}.
    \end{equation}  
    Dễ thấy: $h_1+h_2+h_3<D$ nên hình nón nước không đến bề mặt chỏm băng.
    
    Khối lượng nước đã chảy ra ngoài:
    \begin{equation}
        m'=\frac{h_1+h_2}{h_1+h_2+h_3}m_{tot}=2.7\times10^{10}\si{kg}.
    \end{equation}
    \end{enumerate}
\end{enumerate}
\textbf{Phổ điểm}
    \begin{center}
    \begin{tabular}{|c|p{8cm}|c|}
    \hline
    \multicolumn{1}{|l|}{Phần} & Nội dung & Điểm thành phần \\ 
    \hline
\multirow{1}{*}{(a)} & Thay số và ra được kết quả.
                     & 0.75
                     \\
                     \hline  
\multirow{3}{*}{(b)(i)} & Áp suất tại một điểm dưới lớp băng                                  
                     & 0.25             
                     \\
                     \cline{2-3}
                     & Xét sự dịch chuyển ảo nhỏ và suy ra các phương trình công. 
                     & 0.25
                     \\
                     \cline{2-3}
                     & Áp dụng nguyên lí công áo và suy ra phương trình \eqref{eq:S5_06}. 
                     & 0.25
                     \\
                     \cline{2-3}
                     & Ra được phương trình tương tự phương trình \eqref{eq:S5_10}. 
                     & 0.25
                     \\
                     \cline{2-3}
                     & Kết luận được hình dạng của phần nước. 
                     & 0.25
                     \\
                     \hline
\multirow{4}{*}{(b)(ii)}& Biện luận quá trình diễn ra hiện tượng.
                     & 0.5    
                     \\
                     \cline{2-3}
                     & Xác định độ cao H của đỉnh nón nước dưới chỏm băng. 
                     & 0.5
                     \\
                     \cline{2-3}
                     & Xác định độ cao $h_1$ của magma đã thâm nhập. 
                     & 0.5
                     \\
                     \cline{2-3}
                     & Xác định khối lượng m' của nước. 
                     & 0.5
                     \\
                     \hline
    \end{tabular}
    \end{center}
